% Canonical Paper II source.
The analytic question posed at the beginning of this paper has two answers, because
Paper I supplies two geometrically different carriers for analysis.  A regular AES
is an oriented Riemannian surface whose metric, orientation, assignment, and
arithmetic frame are already fixed.  Its complex analysis is elliptic and intrinsic.
The contact state space is three-dimensional and nonintegrable; its CR and
sub-Riemannian analysis starts only after a horizontal metric and a measure have
been declared.  Keeping these carriers separate is the main organizing decision of
the paper.

On the surface branch, the arithmetic frame makes the first-order drift of the
Laplace--Beltrami operator visible.  If
$[X_u,X_v]=c_uX_u+c_vX_v$, then
\[
 \Delta_g=X_u^2+X_v^2-c_vX_u+c_uX_v.
\]
This formula resolves the ambiguity between a raw sum of squares and the
divergence-form operator selected by Riemannian volume.  The same structure
coefficients give an exact Cauchy--Riemann factorization, and hence a direct proof
that arithmetic holomorphic fields are Laplace--Beltrami harmonic.  The unitary
gauge relation with the ordinary Riemann-surface $\overline\partial$-operator also
fixes the strength of this statement: arithmetic holomorphicity is not a new
holomorphic function class.  It is the classical class expressed in the canonical
frame selected by the assignment law.

On the contact branch, contact volume selects
\[
 \Delta_{\mathcal C}=D_u^2+D_v^2+\lambda D_v,
\]
whereas the raw sum $D_u^2+D_v^2$ is the expression appearing in the elementary
CR product.  The two roles are compatible but not interchangeable.  The adjoint
factorization displays the Reeb correction, while the raw factorization displays
the vertical curvature $\mu\lambda\partial_a$.  The logarithmic first integrals
and affine--Appell modules are consequently explicit CR constructions on the
contact space, not harmonic functions on a hidden horizontal surface.

\subsection{What is established on the basic model}

The complete basic hyperbolic AES provides the global test case.  After the fixed
normalization $X=(\lambda/\mu)x$, it is a constant curvature rescaling of the
standard upper half-plane.  This observation is elementary, but its consequences
for the AEG variables are precise.  The arithmetic coordinate is
\[
 w=X+iy=y\left(i-\frac{\lambda}{\mu}a\right),
\]
the ideal-boundary Poisson kernel solves the $C_0$ Dirichlet problem uniquely, and
the variational conormal derivative has symbol $|D|$.  For Poisson extensions, the
exact energy identity is precisely the homogeneous $H^{1/2}$ boundary quadratic
form and extends by completion from Schwartz data.  Thus the basic model supports
more than a formal differential calculus: it has a proved boundary kernel, a
uniqueness class, a Fourier-multiplier Dirichlet-to-Neumann map on Schwartz data,
and a variational interpretation.

The assignment does not merely relabel these results.  In the $(a,y)$ chart it
turns every admissible holomorphic function of $w$ into an explicitly
assignment-dependent field.  Positive-frequency modes and rational fields supply
concrete families, while the ordinary differential equation for $h(a)$ gives the
complete classification
\[
 h(a)=A+B\arctan\!\left(\frac{\lambda a}{\mu}\right)
\]
of harmonic fields depending only on the assignment.  The nonconstant member is a
boundary half-line harmonic measure.  Its discontinuous trace also shows why the
function class must be stated in every boundary theorem.

\subsection{Holomorphic pullback as an analytic export}

The planar and cylindrical constructions add a second kind of conclusion to the
global boundary theorem.  They show that regular AES data can be transported
functorially by a locally biholomorphic map.  For the planar harmonic base, the
pullback preserves
\[
 \Delta_g a=0,
 \qquad
 K_g=\lambda^2
 \frac{\mu^2-\lambda^2a^2}{\mu^2+\lambda^2a^2}.
\]
For the complete flat cylinder, it preserves
\[
 \Delta_g a=\lambda^2a,
 \qquad
 Z(a)=W^{-1}(S^1).
\]
These formulas give an analytic interface to automorphic constructions: unitary
multipliers preserve the full AES data, while inversion preserves the metric and
zero set and reverses the assignment.  The latter case is naturally expressed by a
sign line bundle on the quotient rather than by silently forcing a scalar-valued
assignment.

This export has a strict boundary.  At a critical point of the holomorphic map the
pullback metric degenerates.  At a zero or pole of the cylindrical target map the
logarithmic assignment ceases to be finite, and an algebraic branch point requires
a choice of local sheet.  Paper II proves no continuation through these sets.  It
supplies the exact formulas on their complement.  Paper III uses this interface to
analyze local branched cones and the specific $q=4$ sign cover; its successful
completion in that model is not a general continuation theorem for other branch
sets.

Table~\ref{tab:final-branch-comparison} summarizes the two analytic branches and the
operator choices proved in this paper.

\begin{table}[t]
\centering
\small
\begin{tabular}{@{}p{0.21\textwidth}p{0.35\textwidth}p{0.35\textwidth}@{}}
\toprule
Feature & Regular-AES surface & Contact state space \\
\midrule
Carrier & oriented Riemannian surface $(M,g)$ & contact three-manifold
$(\mathcal C,\ker\alpha)$ \\
Preferred fields & tangent frame $(X_u,X_v)$ & horizontal lifts
$(D_u,D_v)$ \\
Complex datum & determined by $g$ and orientation & chosen from the normalized
horizontal metric and orientation \\
Natural measure & Riemannian area $d\vol_g$ & positive contact volume
$|\alpha\wedge d\alpha|$ \\
Variational operator & $\Delta_g$, including structure drift &
$\Delta_{\mathcal C}=D_u^2+D_v^2+\lambda D_v$ \\
First-order kernel & ordinary holomorphic class in arithmetic gauge & contact-CR
class with Reeb/vertical twist \\
Global result here & Poisson, Dirichlet, energy, and DtN on the basic model &
Friedrichs realization, hypoellipticity, and explicit CR fields \\
\bottomrule
\end{tabular}
\caption{Analytic data and conclusions on the two carriers.  Similar first-order
notation does not identify the two geometries.}
\label{tab:final-branch-comparison}
\end{table}

\subsection{Analytic frontier}

The results above close the first function-theoretic layer, but they do not close the
programme.  The following problems remain genuinely open within the scope of Paper
II.

\begin{problem}[General boundary theory]
Identify geometric hypotheses on a complete regular AES under which a meaningful
ideal boundary, harmonic measure, Green kernel, and a well-posed Dirichlet class can
be constructed.  The assignment law alone does not provide such a compactification.
\end{problem}

\begin{problem}[Spectral realization beyond the basic model]
Determine the self-adjoint spectral theory of $-\Delta_g$ and of the contact
Friedrichs operator for geometrically controlled nonhomogeneous examples.  In
particular, separate pointwise eigen-equations from $L^2$ eigenvectors and specify
the boundary conditions of every realization.
\end{problem}

\begin{problem}[Contact boundary and completeness]
Develop a boundary-value theory for the normalized contact-CR structure and decide
whether an appropriate completion of the explicit logarithmic and Fourier--CR
families gives a useful representation theorem.  No completeness assertion is made
here.
\end{problem}

\begin{problem}[Degenerate and singular assignments]
Extend the analytic definitions across points where the regular frame ceases to
exist.  This requires the singular geometric input assigned to Paper III; the
present paper neither assumes nor proves a removable-singularity theorem there.
\end{problem}

\begin{problem}[Automorphic singular completion]
Let an automorphic or algebraic function produce a cylindrical pullback on the
regular locus defined by \(0<|W|<\infty\) and \(dW\ne0\).  Determine local and
global conditions under which the metric, assignment, or assignment line bundle
admits a singular AES completion.  Relate ramification indices to zero-graph
valence without assuming that every critical point lies over the unit circle.
\end{problem}

The dependency boundary is therefore explicit.  Paper II uses only the regular
surface, hyperbolic model, and contact connection proved in Paper I.  It exports
elliptic and contact analytic tools, together with planar and cylindrical pullback
formulas, that may later be tested on singular families.  It does not import tube
topology from Paper III or projective-complexity claims from Paper IV.  Within that
boundary, every advertised result is backed by a stated measure, regular locus, the
relevant test-function or variational domain, and a proof.  This is the appropriate
starting point for a broader arithmetic real function theory.
